\documentclass[11pt]{article}
\usepackage[T1]{fontenc}
\usepackage[utf8]{inputenc}
\usepackage{lmodern,microtype}
\usepackage{amsmath,amssymb,amsthm,mathtools}
\usepackage[a4paper,margin=25mm,headheight=15pt]{geometry}
\usepackage{booktabs,array,longtable,enumitem,fancyhdr,xurl,xcolor}
\usepackage[colorlinks=true,allcolors=blue!45!black]{hyperref}
\hypersetup{pdftitle={Cubic-residue quadruplets: elliptic formulas and arithmetic fluctuations},pdfauthor={Brice Pouly}}
\pagestyle{fancy}\fancyhf{}
\fancyhead[L]{\small Cubic-residue quadruplets}
\fancyhead[R]{}
\fancyfoot[C]{\thepage}\renewcommand{\headrulewidth}{0.3pt}
\newtheorem{theorem}{Theorem}[section]
\newtheorem{proposition}[theorem]{Proposition}
\newtheorem{lemma}[theorem]{Lemma}
\newtheorem{corollary}[theorem]{Corollary}
\theoremstyle{definition}\newtheorem{definition}[theorem]{Definition}
\newtheorem{remark}[theorem]{Remark}
\newcommand{\Fp}{\mathbb F_p}
\newcommand{\Q}{\mathbb Q}
\newcommand{\Z}{\mathbb Z}
\newcommand{\E}{\mathbb E}
\newcommand{\Pcal}{\mathcal P}
\newcommand{\Sym}{\operatorname{Sym}}
\newcommand{\Rea}{\operatorname{Re}}
\newcommand{\Nm}{\operatorname{N}}
\newcommand{\ind}{\mathbf1}
\newcolumntype{P}[1]{>{\raggedright\arraybackslash}p{#1}}
\title{Cubic-residue quadruplets:\\
elliptic formulas and arithmetic fluctuations}
\author{Brice Pouly\thanks{Prepared with artificial-intelligence assistance; see the declaration at the end.}}
\date{12 October 2026}
\hypersetup{pdfsubject={Publication edition v1.0; 12 October 2026}}
\begin{document}
\maketitle
\begin{abstract}
Let $A_p$ count the monochromatic consecutive four-windows of a nontrivial
cubic character modulo a prime $p\equiv1\pmod3$, with zero excluded and
overlapping windows counted separately. We give an exact formula for
$A_p$ in terms of fixed elliptic traces and a finite boundary correction.
Its complex-multiplication terms are related cubic twists. This produces
two arithmetic fluctuation regimes: the first four moments of
$(27A_p-p)/\sqrt p$ tend to $(0,44,0,4224)$ when three is a cube and
to $(0,17,0,738)$ otherwise. Over all admissible primes they tend to
$(0,26,0,1900)$. The proof combines a four-position identity following
from Greene's hypergeometric formalism and Lennon's trace formula with
symmetric-power automorphy, cyclic base change and a finite collection of
Rankin--Selberg prime number theorems; it uses no GRH. A direct geometric
proof exposes the phases and exceptional points, with the genus-two and
isogeny constructions of Freeman--Satoh and Guillevic--Vergnaud explicitly
credited. Point-order and Hasse certificates verify numerical examples
without enumerating their billion-position character words. The
conclusions concern fixed pattern length and four moments; no complete
limiting distribution, constant-order centring of the raw count, or law
for logarithmic-length maxima is proved.
\end{abstract}
\noindent\textbf{Keywords.} Cubic characters; Jacobi sums; elliptic
curves; isogenies; complex multiplication; moments over primes;
Rankin--Selberg functions; exact certificates.

\section{The observable and the two main results}
The number $A_p$ of monochromatic cubic four-windows has leading scale
$p/27$, but its fluctuations depend on a small arithmetic condition.
After normalization by $\sqrt p$, the limiting variance is $44$ when
three is a cube modulo $p$, and $17$ otherwise. The cause is not a change
in the number of colours: two complex-multiplication traces in the exact
count are related twists, and their relation depends on the cubic class
of three. We count all overlapping consecutive windows of length four
on the ordinary positions $1,\ldots,p-1$. This observable differs from
maximal-run starts and from a quadratic four-window count.

Questions about consecutive power residues and prescribed character
patterns have a long history, including Lehmer--Lehmer~\cite{lehmer},
the computational theorem on cubic triples of
Lehmer--Lehmer--Mills--Selfridge~\cite{lehmermachine}, and Graham's
result on first occurrences of quadruples~\cite{graham}.
Our observable counts all monochromatic four-windows in one period.
The four-position identity used to compute it is a corollary of
Greene~\cite{greene} and Lennon~\cite{lennon}, with an explicit
fractional linear substitution given below. The associated genus-two
family and isogeny were studied by Freeman--Satoh~\cite{freemansatoh}
and Guillevic--Vergnaud~\cite{guillevicvergnaud}; the latter contains
the short Weierstrass model used here.
The contributions developed here are the assembled exact count with
its boundary corrections, its certified evaluations, and the first
four prime-averaged moments of that count. We retain a direct
geometric proof to expose the phase and contributions at infinity;
we do not claim a new general character-sum or isogeny identity.

\paragraph{Scope and relation to the atlas and certificates.}
The published atlas~\cite{atlas} studies a complete finite population of
power-residue runs below $2^{32}$, selected algebraic windows and separate
sampled ternary diagnostics. These finite computations provide context
for questions about arithmetic structure.
The exact formula and the moment arguments below are self-contained
apart from their cited mathematical inputs. Rabung certificates and
recurrence-closed bounds~\cite{bounds} use finite run maxima and boundary
conditions. Here the observable is instead a count of fixed four-windows,
and its averaging population consists of primes. The finite atlas diagnostics
and the prime-averaged moments below concern different populations. No theorem
about logarithmic-length maxima or new van der Waerden bound is inferred from
the four moments.
When needed, $W(r,k)$ denotes $r$ colours and progression length $k$;
the ordinary window convention is fixed explicitly below.

\subsection{Windows, characters, and the prime population}
Throughout, $p>3$ is prime, $p\equiv1\pmod3$, and $\chi$ is a
nontrivial cubic character of $\Fp^\times$. Write
$\omega^2+\omega+1=0$ and regard its values as
$1,\omega,\omega^2$. Extend $\chi$ by $\chi(0)=0$.
Because $3\mid(p-1)/2$, we have $\chi(-1)=1$.
Conjugating $\chi$ does not change any equality of colours.

Define the cubic four-window count by
\begin{equation}\label{eq:count}
 A_p=\#\{1\le x\le p-4:
       \chi(x)=\chi(x+1)=\chi(x+2)=\chi(x+3)\}.
\end{equation}
All three nonzero colours are included. The four entries are consecutive
integer residues, and overlapping windows are counted separately.
This is a window statistic, not a count of maximal runs. The character
order is part of its definition: a cubic four-window count is not
reconstructed by relabelling a quadratic four-window count.

Put
\begin{equation}\label{eq:correction}
 I_2(p)=\ind_{\chi(2)=1},\qquad
 I_{23}(p)=\ind_{\chi(2)=\chi(3)=1},\qquad
 b_p=26+18I_2(p)+18I_{23}(p).
\end{equation}
Thus $b_p\in\{26,44,62\}$. We use both normalizations
\begin{equation}\label{eq:DY}
 D_p=\frac{27A_p-p+b_p}{\sqrt p},
 \qquad Y_p=\frac{27A_p-p}{\sqrt p}.
\end{equation}
The first retains the exact finite boundary correction; the second
is the more immediate fluctuation around $p/27$.

Our population of varying primes is
\[
 \Pcal(X)=\{p\text{ prime}:X\le p<2X,\ p\equiv1\pmod3\}.
\]
Let $\Pcal_0(X)$ consist of its primes for which three is a cube,
and let $\Pcal_*(X)$ be the complement. Each average below is
uniform on the stated finite set and is divided by that set's own
cardinality. There is no additional random origin in a character
word, or additional random choice of the orientation of $\chi$.

\subsection{An exact elliptic formula}
For an elliptic curve $E/\Fp$ write
$t_E(p)=p+1-\#E(\Fp)$. Consider the fixed rational curves
\begin{align}
 E_1&:y^2=x^3-99x+414,\nonumber\\
 E_2&:y^2=x^3-21x+37,\label{eq:noncm}\\
 E_3&:y^2=x^3-2835x+58158,\nonumber
\end{align}
and the three curves with complex multiplication
\begin{equation}\label{eq:cmcurves}
 C_0:y^2=x^3-27,\quad
 C_1:y^2=x^3-243,\quad
 C_2:y^2=x^3-2187.
\end{equation}
All six have good reduction for $p>3$. Write $t_i=t_{E_i}$,
$c_j=t_{C_j}$, and
$\beta_p=\chi(3)+\overline{\chi(3)}\in\{2,-1\}$.

\begin{theorem}[Exact four-window formula]\label{thm:exact}
For every admissible prime,
\begin{equation}\label{eq:exact}
 27A_p=p-b_p-2c_0(p)-2c_1(p)
                    -2t_1(p)-2t_2(p)-\beta_p t_3(p).
\end{equation}
Equivalently, if $J_p=J(\chi,\chi)$,
\begin{align}
 27A_p={}&p-b_p+
 4\Rea\{J_p(\overline{\chi(2)}+\overline{\chi(6)})\}
 \nonumber\\[-2pt]
 &-2t_1(p)-2t_2(p)-\beta_p t_3(p).
 \label{eq:jacobi-form}
\end{align}
The CM traces satisfy
\begin{equation}\label{eq:twistrelations}
 \begin{cases}
 c_0=c_1=c_2,& \chi(3)=1,\\
 c_0+c_1+c_2=0,& \chi(3)\ne1.
 \end{cases}
\end{equation}
\end{theorem}
Sections~\ref{sec:fourpositions}--\ref{sec:CM} prove these formulas.
The finite corrections in \eqref{eq:exact} matter for an exact count:
the character expansion on all of $\Fp$ is not the desired count
until the four excluded residues have been removed.

\subsection{Four moments over primes}
\begin{theorem}[Moments over the actual primes]\label{thm:moments}
As $X\to\infty$, the relative cardinalities of $\Pcal_0(X)$ and
$\Pcal_*(X)$ inside $\Pcal(X)$ tend to $1/3$ and $2/3$.
For $r=1,2,3,4$, the limits of the averages of $D_p^r$ are
\[
 \begin{array}{c|rrrr}
 \text{population}&r=1&r=2&r=3&r=4\\ \hline
 \Pcal_0(X)&0&44&0&4224\\
 \Pcal_*(X)&0&17&0&738\\
 \Pcal(X)&0&26&0&1900
 \end{array}
\]
The same statements hold with $Y_p$ in place of $D_p$.
\end{theorem}
The constants can be read from the trace combination before the analytic
proof. Write $x_i=t_i/\sqrt p$ and $v_j=c_j/\sqrt p$. Equations
\eqref{eq:exact} and \eqref{eq:twistrelations} give
\[
 D_p=\begin{cases}
 -4v_0-2x_1-2x_2-2x_3,&\chi(3)=1,\\
 2v_2-2x_1-2x_2+x_3,&\chi(3)\ne1.
 \end{cases}
\]
The low-degree averaging rule proved below assigns second moment $2$ to
the CM coordinate and $1$ to each non-CM coordinate, with vanishing mixed
terms in the required range. Thus
\[
 44=4^2\cdot2+2^2+2^2+2^2,\qquad
 17=2^2\cdot2+2^2+2^2+1.
\]
Proving this averaging rule under the finite cubic conditioning is the
analytic task. It provides the needed mixed moments through total degree
four and does not posit an independent joint law in every degree.

The analytic proof in Sections~\ref{sec:qualification}--\ref{sec:moments}
uses fixed automorphic objects. It establishes the mixed moments
needed up to total degree four, without postulating complete joint
equidistribution of the six curves.
The displayed limits have $o(1)$ errors; we do not extract a
numerical error bound applicable near $2^{32}$.

\section{A four-position identity from Greene and Lennon}\label{sec:fourpositions}
\subsection{The phase and the elliptic curve}
Take distinct $a,b,c,d\in\Fp$ and set
\begin{equation}\label{eq:parameters}
 \rho=\frac{c-d}{a-b},\qquad
 A=\frac{4(ab+cd)-2(a+b)(c+d)}{(a-b)(c-d)}.
\end{equation}
Define
\begin{equation}\label{eq:EA}
 E_A:Y^2=X^3-(6A+15)X+(A^2+14A+22),\qquad
 t_A(p)=p+1-\#E_A(\Fp).
\end{equation}
The identities
\[
 A-2=\frac{4(a-d)(b-c)}{(a-b)(c-d)},\qquad
 A+2=\frac{4(a-c)(b-d)}{(a-b)(c-d)}
\]
give $A\ne\pm2$. Direct substitution in the Weierstrass
discriminant gives
\begin{equation}\label{eq:disc}
 \Delta(E_A)=-432(A-2)^3(A+2)\ne0.
\end{equation}
Thus none of the curves occurring in the identity is singular.

\begin{corollary}[Four-position correlation; Greene--Lennon]\label{thm:fourpositions}
With the conventions above,
\begin{equation}\label{eq:fouridentity}
 1+\sum_{z\in\Fp}
 \chi((z+a)(z+b)(z+c)^2(z+d)^2)
 =-\overline{\chi(\rho)}\,t_A(p).
\end{equation}
Moreover $t_A(p)=t_{-A}(p)$.
If $\beta_p(u)=\chi(u)+\overline{\chi(u)}$, the conjugate pair
of sums is therefore
\begin{equation}\label{eq:pair}
 S_e(p)+S_{-e}(p)=-2-\beta_p(\rho)t_A(p).
\end{equation}
\end{corollary}
The explicit phase is part of the formula. Isomorphism over an
algebraic closure alone would not recover it.

\subsection{The hypergeometric derivation and its normalization}
Extend the trivial character $\varepsilon$ by $\varepsilon(0)=0$,
as in Greene. Put
\[
 s=\frac{2-A}{4}
   =\frac{(a-d)(c-b)}{(a-b)(c-d)},\qquad
 u=\frac{a-b}{a-d}\frac{z+d}{z+b}.
\]
The distinct-position assumption gives $s\ne0,1$. This projective
change of variable sends $-d,-a,-c,-b$ to $0,1,1/s,\infty$, and
\[
 \frac{(z+a)(z+b)}{(z+c)(z+d)}
   =\rho^{-1}\frac{1-u}{u(1-su)}.
\]
Define Greene's normalized function, here specialized to cubic
characters, by its character sum:
\begin{equation}\label{eq:greene}
 F(s)={}_2F_1\!\left(\begin{matrix}\chi,&\overline\chi\\&\varepsilon\end{matrix};s\right)_p
 =\frac1p\sum_{u\in\Fp}
      \overline\chi(u)\chi(1-u)\overline\chi(1-su).
\end{equation}
This is Definition 3.5 of~\cite[p.~82]{greene}, since
$\chi(-1)=1$ and $\varepsilon(s)=1$. The poles and zeros contribute
zero; the source point $z=\infty$ contributes one. Therefore the
left side of \eqref{eq:fouridentity} is
$\overline\chi(\rho)\,pF(s)$, including its cubic phase and the
extra one.

In \eqref{eq:EA} the substitution
$X=x+3$, $Y=y-3x+A-2$ gives
\[
 y^2-6xy+2(A-2)y=x^3.
\]
Lennon's Theorem 1.1~\cite{lennon}, with
$a_1=-6$, $a_3=2(A-2)$, and $27a_3/a_1^3=s$, gives
$t_A(p)=-pF(s)$. A residue value of $A$ can be lifted to an integer;
$p>3$ and $A\ne\pm2$ give the required good reduction and
$p\nmid a_1$. This also covers $A=0$.
Greene's transformation~\cite[Theorem 4.4(i), p.~88]{greene}
specializes to $F(s)=F(1-s)$, so $t_A=t_{-A}$.
Conjugation gives \eqref{eq:pair}. This proves the corollary using
the published identities.

\subsection{Direct geometric proof: a degree-two rational map}
Let
$R(z)=(z+a)(z+b)/((z+c)(z+d))$.
Its character equals that of the polynomial in
\eqref{eq:fouridentity} wherever the quotient is defined, because
the two differ by a cube. Its poles contribute zero, while
$R(\infty)=1$. If $S_e$ denotes the affine polynomial sum,
the character sum on this projective source is $S_e+1$.

For $v\in\Fp$, the homogeneous equation $R(z)=v$ has
$1+\lambda(D(v))$ distinct rational roots, where $\lambda$ is
the quadratic character and
\[
 D(v)=(c-d)^2v^2+
 [4(ab+cd)-2(a+b)(c+d)]v+(a-b)^2.
\]
This root count includes infinity when the affine leading
coefficient vanishes. The equation is never identically zero:
its numerator and denominator have four distinct roots.
Since $\sum_v\chi(v)=0$, changing variables by $w=\rho v$ gives
\begin{equation}\label{eq:mixed}
 S_e+1=\overline{\chi(\rho)}B_A,\qquad
 B_A=\sum_{w\in\Fp}\chi(w)\lambda(w^2+Aw+1).
\end{equation}
Inversion $w\mapsto w^{-1}$, with the zero term treated separately,
shows that $B_A=\overline{B_A}$. Since it belongs to $\Z[\omega]$,
it is an integer. We will identify it with $-t_A$.

\subsection{The genus-two curve and its elliptic quotients}
Let $C_A$ be the smooth completion of
$y^2=x^6+Ax^3+1$. The polynomial is squarefree because $A^2\ne4$,
and the leading coefficient gives two rational points at infinity.
This is the $b=1$ specialization of the family in
Freeman--Satoh~\cite[Proposition 4.2]{freemansatoh}.
The short model \eqref{eq:EA} and its degree-three isogeny also
appear in Guillevic--Vergnaud~\cite[Section 3.1, p.~241]{guillevicvergnaud}.
We give the maps and an elementary point count for this specialization.
The number of cube roots of $w$ is
$1+\chi(w)+\overline{\chi(w)}$, including the value one at zero.
Using $\sum_w\lambda(w^2+Aw+1)=-1$ gives
\[
 \sum_x\lambda(x^6+Ax^3+1)=-1+2B_A,\qquad
 \#C_A(\Fp)=p+1+2B_A.
\]
Thus its genus-two trace, defined as $p+1-\#C_A(\Fp)$,
is $-2B_A$.

The two quotient curves are
\[
 E_+:v^2=(u+2)(u^3-3u+A),\qquad
 E_-:v^2=(u-2)(u^3-3u+A).
\]
On suitable affine opens the maps are
$u=x+x^{-1}$ and $v=y(x\pm1)/x^2$.
We keep track of the traces by an elementary count, rather than
requiring a decomposition theorem for the Jacobian.
For each $u$,
\begin{equation}\label{eq:fibers}
 \sum_{\substack{x\ne0\\x+x^{-1}=u}}\lambda(x)
 =\lambda(u+2)+\lambda(u-2).
\end{equation}
Indeed, away from $u=\pm2$, the two roots have the same
quadratic character, and $(x\pm1)^2/x=u\pm2$. A nonsplit fibre
has zero on both sides in aggregate. At $u=\pm2$ the formula
is checked with the single root $x=\pm1$.
Now
$x^6+Ax^3+1=x^3(u^3-3u+A)$ for $x\ne0$.
Equation~\eqref{eq:fibers}, the term $x=0$, and the two
infinities on each quartic yield
\begin{equation}\label{eq:trace-split}
 t(C_A)=t(E_+)+t(E_-).
\end{equation}

The substitutions
\[
 w=(u+2)^{-1},\quad v'=vw^2,\quad
 X=(A-2)w+3,\quad Y=(A-2)v'
\]
transform $E_+$ to \eqref{eq:EA}.
Birational maps between their smooth projective completions
preserve the point count, including the points outside this open.
The change $u\mapsto-u$ identifies $E_-$ with the corresponding
quartic for $-A$, hence with $E_{-A}$.

\subsection{The degree-three isogeny and the square twist}
The point $Q=(3,A-2)$ on $E_A$ has order exactly three.
Its tangent has slope $-3$, and the group law gives $2Q=-Q$.
With $h=X-3$, put
\[
 T=-12(A-2),\qquad V=4(A-2)^2,\qquad
 r(X)=X+\frac{T}{h}+\frac{V}{h^2}.
\]
The map $(X,Y)\mapsto(r(X),r'(X)Y)$ is the
V\'elu isogeny with kernel $\{O,Q,-Q\}$, whose target is
\begin{equation}\label{eq:velu-target}
 Y^2=X^3+(54A-135)X-27A^2+378A-594.
\end{equation}
The rational identity defining this map can also be verified
directly over $\Z[A,X]$ after multiplication by $h^6$;
the supplied verifier does this formal calculation.
Since $p>3$ the isogeny is separable and extends across its
kernel to the smooth projective curves.

The target coefficients are those of the quadratic twist of
$E_{-A}$ by $-3$: they are nine times its linear coefficient
and $-27$ times its constant coefficient. A primitive cubic
root $\zeta\in\Fp$ satisfies $(2\zeta+1)^2=-3$.
Consequently this twist is isomorphic to $E_{-A}$ over $\Fp$.
Curves isogenous over a finite field have the same trace.
For this direction, an isogeny commutes with Frobenius minus
the identity; taking degrees and cancelling its degree proves
equality of the degrees of Frobenius minus the identity,
which are the point counts.

We have proved $t(E_+)=t(E_-)=t_A(p)$.
Combining this with \eqref{eq:trace-split} gives
$-2B_A=2t_A(p)$. Equation~\eqref{eq:mixed} proves
\eqref{eq:fouridentity}, and conjugation proves \eqref{eq:pair}.
This completes the direct proof of Corollary~\ref{thm:fourpositions}.
The two elliptic quotients need not be isomorphic to each other;
their equality of traces here follows from the isogeny and the
square twist.

\section{Derivation of the consecutive quadruplet formula}
\subsection{The pair and triple sectors}
Away from the four zeros, the equality indicator has the expansion
\begin{equation}\label{eq:expansion}
 27\,\ind_{\chi(z)=\cdots=\chi(z+3)}
 =\sum_{\substack{e\in\{0,1,2\}^4\\\sum e_i=0\ (3)}}
   \prod_{i:e_i\ne0}\chi(z+i)^{e_i}.
\end{equation}
An empty product is one. This convention for zero exponents
is retained at the excluded residues when the polynomial
expansion is summed over $\Fp$.
There are 27 terms: the constant, 12 of support two, eight
of support three, and six of support four.

Every support-two sum is $-1$. For example, the quotient
$(z+a)/(z+b)$ parametrizes $\Fp\setminus\{1\}$ as $z$ runs
away from its pole, and has the same cubic character as
$(z+a)(z+b)^2$. Its full sum is therefore $-1$.

For distinct $i<j<l$, set
$\delta=(j-i)(l-i)(l-j)$.
We claim the exact triple identity
\begin{equation}\label{eq:triple}
 1+\sum_{z\in\Fp}\chi((z+i)(z+j)(z+l))
 =J_p\overline{\chi(\delta)}.
\end{equation}
To see the missing constant explicitly, let
$t=(l-j)(z+i)/((l-i)(z+j))$.
It sends the three roots to $0,\infty,1$ and satisfies
\[
 (z+i)(z+j)(z+l)=
 \delta^{-1}\big((l-i)(z+j)\big)^3t(1-t).
\]
The value of $t$ omitted by finite $z$ is
$t_\infty=(l-j)/(l-i)$.
Since $\chi(t_\infty(1-t_\infty))=\chi(\delta)$,
the Jacobi sum loses exactly one after multiplication by
$\overline{\chi(\delta)}$. The pole at $z=-j$ contributes zero.
This proves \eqref{eq:triple} with its orientation and its constant.

The four triples of $\{0,1,2,3\}$ have difference products
$2,6,6,2$. Adding their conjugates gives the nonconstant
contribution
$4\Rea\{J_p(\overline{\chi(2)}+\overline{\chi(6)})\}-8$.

\subsection{The four-support sector and all four exclusions}
For the three unoriented partitions of $\{0,1,2,3\}$, the
parameters $(A,\rho)$ in \eqref{eq:parameters} are
\[
 (14,1),\qquad(-1,1),\qquad(-10/3,1/3).
\]
Their conjugate sums contribute
$-6-2t_{14}-2t_1-\beta_p t_{10/3}$.
The equalities of opposite-parameter traces have been proved,
and $\beta_p(1/3)=\beta_p(3)$.
The first two positive parameters give $E_1,E_2$ in
\eqref{eq:noncm}. The map
$(X,Y)\mapsto(x,y)=(9X,27Y)$ sends $E_{10/3}$,
whose constant coefficient is $718/9$, to $E_3$.
It is an isomorphism for every $p>3$.
Together with the constant and pair sectors, the complete
sum over $\Fp$ has main term $p-26$.

We must remove $z=-r$, $r=0,1,2,3$.
Terms with an active exponent at $r$ are already zero.
The remaining terms sum to nine if the other three values
are equal, and to zero otherwise. At $r=0,3$ the equality
condition is $\chi(2)=\chi(3)=1$; at $r=1,2$ it is
$\chi(2)=1$, using $\chi(-1)=1$.
Thus the exclusions subtract exactly
$18I_2+18I_{23}$, including their share of the constant term.
The remaining residues are precisely $z=1,\ldots,p-4$,
and \eqref{eq:jacobi-form} follows.
This calculation accounts for all zero terms and avoids an
error of four from treating the full field as the nonzero interval.

\section{Complex multiplication and the exact conditional form}\label{sec:CM}
\subsection{Fermat curves and Jacobi traces}
For $a\in\Fp^\times$, the smooth projective Fermat cubic
$u^3+v^3=aw^3$ has three rational points at infinity.
Counting cube fibres in the affine chart gives
\begin{equation}\label{eq:fermat-count}
 \#F_a(\Fp)=p+1+
 J_p\overline{\chi(a)}+\overline{J_p}\chi(a).
\end{equation}
Indeed the two mixed character terms each equal $-1$;
the remaining nonconstant pair is the two Jacobi terms.
The affine count is $p-2$ plus that pair, and the three
infinities add three. This proves the sign in
\eqref{eq:fermat-count} directly.

On $u+v\ne0$, the change
\[
 X=\frac{12a}{u+v},\qquad
 Y=\frac{36a(u-v)}{u+v}
\]
gives $Y^2=X^3-432a^2$.
Conversely $u+v=12a/X$ and $u-v=Y/(3X)$.
These birational maps identify the smooth projective curves,
so the exceptional points do not alter the point count.
For $a=2,6,18$, scaling $(X,Y)=(4x,8y)$ gives
the three curves in \eqref{eq:cmcurves}. Hence
\begin{equation}\label{eq:cm-jacobi}
 c_j(p)=-2\Rea\{J_p\overline{\chi(2\cdot3^j)}\},
 \qquad j=0,1,2.
\end{equation}
Substitution proves \eqref{eq:exact}.
If $\eta=\chi(3)$ is one, the three traces coincide.
Otherwise $1+\eta+\eta^2=0$ proves their zero sum.
This also proves \eqref{eq:twistrelations}.

\subsection{Two arithmetic fluctuation regimes}
Put $x_i=t_i/\sqrt p$ and $v_j=c_j/\sqrt p$.
Equations~\eqref{eq:exact} and \eqref{eq:twistrelations} give
\begin{equation}\label{eq:conditional}
 D_p=
 \begin{cases}
 -4v_0-2x_1-2x_2-2x_3,&3\text{ a cube},\\
  2v_2-2x_1-2x_2+x_3,&3\text{ not a cube}.
 \end{cases}
\end{equation}
Hasse gives $|v_j|,|x_i|\le2$. Thus $|D_p|\le20$
in the first class and $|D_p|\le14$ in the second.
Since $|D_p-Y_p|\le62/\sqrt X$ on $[X,2X)$,
the limits of every fixed moment agree for these two
normalizations whenever either limit is established.

The two CM terms in the original exact formula are related
twists. Treating them as independent would give the wrong
second and fourth moments. Their dependence also persists
when their unconditional covariance happens to vanish.

\section{Arithmetic qualification of the non-CM factors}\label{sec:qualification}
\subsection{Non-CM and good reduction}
The $j$-invariants of $E_1,E_2,E_3$ are, respectively,
$-35937/4$, $790272$, and $-3472875/4$.
The first and third are not algebraic integers, which
already excludes complex multiplication for them.
A uniform alternative for all three curves uses two
good ordinary reductions with different Frobenius fields:
\begin{center}
\begin{tabular}{c rr l rr l}
\toprule
Curve&$p$&$t_p$&Frobenius field&$p$&$t_p$&Frobenius field\\
\midrule
$E_1$&5&$-3$&$\Q(\sqrt{-11})$&7&$-4$&$\Q(\sqrt{-3})$\\
$E_2$&5&3&$\Q(\sqrt{-11})$&7&$-1$&$\Q(\sqrt{-3})$\\
$E_3$&7&$-1$&$\Q(\sqrt{-3})$&19&5&$\Q(\sqrt{-51})$\\
\bottomrule
\end{tabular}
\end{center}
In each entry $p\nmid t_p$, so the reduction is ordinary.
The specialization of endomorphisms is injective at good
reduction~\cite{sutherland21}. For an ordinary reduction,
the geometric endomorphism algebra is the imaginary
quadratic field $\Q(\sqrt{t_p^2-4p})$~\cite{sutherland13}.
This remains the same on finite extensions: the two
Frobenius roots have different $p$-adic valuations, so their
ratio is not a root of unity and no power of Frobenius
becomes rational. A characteristic-zero CM field would
therefore inject into each of the two different fields
in its row, which is impossible.

The Weierstrass discriminants are
$-2^{14}3^6$, $2^4 3^4$, and $-2^{14}3^{11}$.
Thus all primes under consideration are good.
The trace witnesses and these elementary invariants are
recomputed in the finite verifier.

\subsection{Symmetric powers and quadratic base change}
Let $\pi_i$ denote the unitary-normalized modular
representation attached to $E_i/\Q$. Its central character
is trivial and it is self-contragredient.
Modularity and Newton--Thorne's Theorem A~\cite{newtonthorne}
provide cuspidal $\Sym^m\pi_i$ on $\mathrm{GL}_{m+1}/\Q$
for $1\le m\le4$.
Set $K=\Q(\omega)$. The cyclic base-change criterion
recalled in~\cite[Proposition 2.3.1(5), p.~54]{ramakrishnan} says that a
cuspidal representation remains cuspidal in this quadratic
extension unless it has the corresponding quadratic
self-twist.

This exceptional self-twist is excluded by a single inert
prime for each curve and every needed $m$.
Define $B_0=1$, $B_1=t$, and
$B_m=tB_{m-1}-pB_{m-2}$.
The normalized symmetric-power trace is $B_m/p^{m/2}$.
At the following inert primes, none is zero:
\begin{center}
\begin{tabular}{c r rrrr}
\toprule
Curve&$p$&$B_1$&$B_2$&$B_3$&$B_4$\\
\midrule
$E_1$&5&$-3$&4&3&$-29$\\
$E_2$&5&3&4&$-3$&$-29$\\
$E_3$&11&$-6$&25&$-84$&229\\
\bottomrule
\end{tabular}
\end{center}
A self-twist by the quadratic character of $K/\Q$
would negate each such trace, forcing it to be zero.
Write $\Pi_i$ for the base change of $\pi_i$ to $K$.
The representations $\Sym^m\Pi_i$ are consequently
cuspidal, unitary, and self-contragredient, with trivial
central character. In particular $\Pi_i$ is nondihedral,
because a dihedral $\mathrm{GL}_2$ representation has a
noncuspidal symmetric square.

\subsection{Twist exclusions and the tensor products}
At either place of $K$ above the split prime 31 the traces
are
\begin{equation}\label{eq:31}
 (t_1,t_2,t_3)=(-4,5,-7).
\end{equation}
If two $\Pi_i$ were character twists of each other, their
trivial central characters would force the twisting
character to have square one. Both local representations
are unramified at 31, so that character is unramified there
and the absolute values of their traces would agree.
The three distinct values in \eqref{eq:31} exclude this.
Ramakrishnan's tensor-product theorem~\cite[Theorem M]{ramakrishnan}
therefore gives cuspidal, self-contragredient
$\Pi_i\boxtimes\Pi_j$ on $\mathrm{GL}_4/K$ for $i\ne j$.
Here we use both nondihedrality and the twist exclusion.

The normalized symmetric-square traces at 31 are
\begin{equation}\label{eq:sym2-31}
 -15/31,\qquad-6/31,\qquad18/31.
\end{equation}
Their absolute values are also distinct. In the moment
proof all Hecke twists are unramified at 31 and unitary.
Consequently \eqref{eq:31} and \eqref{eq:sym2-31}
exclude the equal-rank collisions even after an arbitrary
imaginary power of the norm is included: such a twist
still multiplies a trace at 31 by a number of modulus one.
This local argument deals with the norm twists appearing
in the prime number theorem as well as the finite cubic
characters. We do not infer independence just from
different equations or different $j$-invariants.

\section{A twisted averaging lemma through degree four}\label{sec:analytic}
\subsection{The CM Hecke character and the finite prescription}
CM Hecke theory~\cite[Theorem 3.6 in arXiv v1]{chenshen} associates to
$C_0$ a unitary Hecke character $\gamma$ of $K$ with
nonzero angular infinity type. At a good prime ideal
$\mathfrak p$, its normalized trace is
$\gamma(\mathfrak p)+\gamma(\mathfrak p)^{-1}$.
The field $K(\sqrt[3]{3})/K$ is cyclic of degree three.
Indeed the valuation of three at the place above three
is two, not a multiple of three.
Let $\eta$ denote its nontrivial cubic Hecke character.
At a degree-one place it records the cubic class of
three, with conjugation coherent with the character
orientation.

The three CM twists may be labelled so that
\begin{equation}\label{eq:gamma}
 \begin{split}
 v_0&=\gamma+\gamma^{-1},\\
 v_1&=\gamma\eta^{-1}+\gamma^{-1}\eta,\\
 v_2&=\gamma\eta+\gamma^{-1}\eta^{-1}.
 \end{split}
\end{equation}
Indeed, if $v^3=3$, the maps
$(x,y)\mapsto(v^{2j}x,v^{3j}y)$ identify $C_0$ with $C_j$
over $K(v)$ for $j=1,2$. Their descent characters are the two
nontrivial powers of $\eta$, giving \eqref{eq:gamma} after choosing
an orientation. Conditioning the CM angle on a finite Galois class
is an established method; see Fité--Sutherland~\cite[Propositions
3.6 and 3.10]{fitesutherland}. The particular six-curve observable
still requires the mixed products below.
Reversing the reciprocity convention conjugates both
characters and leaves the class-conditioned moments
unchanged. Their conductors have support only above
two and three, so all twists here are unramified at 31.

At any good place of $K$ with residue cardinality $q$,
let $x_i$ be the trace of $E_i$ divided by $\sqrt q$.
Define $P_0(x)=1$, $P_1(x)=x$,
$P_m(x)=xP_{m-1}(x)-P_{m-2}(x)$.
Then $P_m(x_i)$ is the local trace of $\Sym^m\Pi_i$.

\begin{lemma}[Fixed-object twisted averages]\label{lem:averages}
For $h\in\Z$, $u\in\{0,1,2\}$, and nonnegative
$m_1,m_2,m_3$ with
$|h|+m_1+m_2+m_3\le4$, the average over good prime
ideals $\mathfrak p$ of $K$ with $X\le\Nm\mathfrak p<2X$ of
\begin{equation}\label{eq:monomial}
 \gamma(\mathfrak p)^h\eta(\mathfrak p)^u
 \prod_{i=1}^3P_{m_i}(x_i(\mathfrak p))
\end{equation}
tends to zero unless $h=m_1=m_2=m_3=u=0$.
In that constant case the average is one.
\end{lemma}

\subsection{The finite list of analytic products}
Let $\theta=\gamma^h\eta^u$.
Gillespie--Ji's Theorem 1.1 in arXiv v1~\cite{gillespieji} applies over
the fixed Galois field $K$ to two unitary cuspidal factors
when at least one is self-contragredient.
Its possible main term comes from equivalence of the
factors up to an imaginary norm twist. We use the
contragredient on the second factor when converting
the notation to a product of the desired local traces.

The required products are exhausted by the following list.
The regrouping of the degree pattern $(2,1,1)$ into ranks $(3,4)$
uses the same tensor-product strategy as
Ramakrishnan~\cite[Theorem 4.2.1, pp.~93--94]{ramakrishnan};
the finite enumeration below checks its applicability to our
specific moment expansion, rather than introducing that strategy.
The indices $i,j,l$ are distinct whenever they occur together.
\begin{center}\small
\begin{tabular}{P{20mm}P{40mm}P{78mm}}
\toprule
Positive $m_i$&Allowed pattern&Factors for the prime number theorem\\
\midrule
None&$\theta$ alone&$\theta$ and the trivial $\mathrm{GL}_1$ character\\
One&$1\le m_i\le4$&$\Sym^{m_i}\Pi_i$ and $\theta$\\
Two&$m_i+m_j\le4$&$\Sym^{m_i}\Pi_i$ and $\Sym^{m_j}\Pi_j\otimes\theta$\\
Three&$(1,1,1)$, $|h|\le1$&
 $\Pi_1\boxtimes\Pi_2$ and $\Pi_3\otimes\theta$\\
Three&Permutation of $(2,1,1)$, $h=0$&
 $\Sym^2\Pi_i$ and $(\Pi_j\boxtimes\Pi_l)\otimes\eta^u$\\
\bottomrule
\end{tabular}
\end{center}
Every listed first factor, except for the first row
where the trivial factor can be used, is cuspidal
and self-contragredient by Section~\ref{sec:qualification}.
If the two ranks differ, equivalence is impossible.
The two three-index rows have ranks $(4,2)$ and
$(3,4)$, respectively. In the two-index row the only
equal-rank patterns are $(1,1)$ and $(2,2)$.
Equations~\eqref{eq:31} and \eqref{eq:sym2-31}
exclude both after every twist under consideration,
including the imaginary norm twist.
The one-index row has ranks $m_i+1$ and one.

In the first row $\theta$ can equal an imaginary norm
power only when $h=0$, by its angular infinity type.
Then its order-three part is nontrivial unless $u=0$.
In the constant case the usual prime ideal theorem
provides the main term. Thus every nonconstant
case has no main term.

\subsection{Prime powers, ramified places, and unweighted averages}
The cited theorem initially controls coefficients of a
logarithmic derivative, with norm logarithms as weights.
We spell out the passage to prime ideals used here.
At every good place, normalized elliptic eigenvalues
have modulus one. Symmetric powers and the two-factor
tensor products retain that property; the Hecke
characters are also unitary. Their traces at prime
powers are therefore bounded by a constant depending
only on the fixed ranks.
The contribution from powers of order at least two up
to norm $2X$ is $O(\sqrt X(\log X)^2)$.

There are finitely many bad places. Their factors here
come from the local parameters of elliptic curves
and CM characters, followed by the same fixed tensor
and symmetric-power operations. In unitary
normalization their local parameters in the Euler
factors have modulus at most one. More explicitly, an elliptic
unitary Weil--Deligne parameter is tempered: its special blocks
have the form $U\otimes\operatorname{Sp}(r)$ with unitary $U$.
On inertia invariants in the kernel of monodromy, the eigenvalues
from such a block have modulus $q^{-(r-1)/2}\le1$.
The fixed symmetric powers and tensor products here decompose into
blocks of this same form; the unitary CM characters have $r=1$.
Thus this bound applies to the actual ramified Rankin--Selberg
Euler factors, not just to the good places.
Their higher-power contribution is consequently
$O(\log X)$ after summing over the fixed bad places.
Equivalently these factors may be removed before
applying the prime number theorem and their finite
Euler corrections restored afterwards.

The nonconstant logarithmically weighted sum is $o(X)$;
the constant sum is $X+o(X)$.
Partial summation on $[X,2X)$ removes the logarithmic
weight and gives $o(X/\log X)$ in the nonconstant cases,
against $X/\log X+o(X/\log X)$ for the prime count.
This proves Lemma~\ref{lem:averages}.
We have paid for prime powers directly using known
tempered parameters. No general Ramanujan conjecture,
Hypothesis H, or growing-conductor uniformity is being
imported from a more general statement.

\section{Conditioning and the integer moment calculations}\label{sec:moments}
\subsection{From prime ideals to rational primes}
Prime ideals of residue degree at least two and norm
at most $2X$ lie over rational primes at most $\sqrt{2X}$.
There are $O(\sqrt X)$ of them in this quadratic field.
All observables needed here are bounded, so their
contribution is negligible in the prime ideal averages.
Each split rational prime $p\equiv1\pmod3$ has exactly
two conjugate degree-one places of norm $p$.
The actual count, the trace combination $D_p$, and the
events $\eta=1$ and $\eta\ne1$ are invariant under
conjugation. Averaging them over the two places is
therefore exactly the rational-prime average in
Theorem~\ref{thm:moments}.

The class indicators are
\begin{equation}\label{eq:indicators}
 \ind_{\eta=1}=\frac{1+\eta+\eta^2}{3},\qquad
 \ind_{\eta\ne1}=\frac{2-\eta-\eta^2}{3}.
\end{equation}
Lemma~\ref{lem:averages} gives relative densities
$1/3$ and $2/3$. Dividing the corresponding
weighted averages by those densities gives the
class-conditioned moments, not moments divided
by the total population.

\subsection{The mixed functional only through degree four}
For nonnegative $r_0,r_1,r_2,r_3$ of total degree
at most four, the conditional averages of
$v_0^{r_0}\prod_i x_i^{r_i}$ in the cube class,
and of $v_2^{r_0}\prod_i x_i^{r_i}$ in the
noncube class, factor as
\begin{equation}\label{eq:functional}
 M_{\rm CM}(r_0)\prod_{i=1}^3M_{\rm ST}(r_i).
\end{equation}
Here, for orders zero through four,
\[
 M_{\rm CM}=(1,0,2,0,6),\qquad
 M_{\rm ST}=(1,0,1,0,2).
\]
To prove this formula, expand \eqref{eq:gamma} in
powers of $\gamma$ and $\eta$, use
\[
 x^2=P_2(x)+1,\quad
 x^3=P_3(x)+2P_1(x),\quad
 x^4=P_4(x)+3P_2(x)+2,
\]
and multiply by \eqref{eq:indicators}.
Every nonconstant term vanishes by
Lemma~\ref{lem:averages}. The surviving constant
coefficients are $\binom{2r}{r}$ for a CM even
power $2r$ and the Catalan number
$\binom{2r}{r}/(r+1)$ for each non-CM factor.
Odd powers have no constant coefficient.
The factors of $\eta$ in $v_2$ cancel in its
constant term, so the same functional applies
in the noncube class.

This factorization is proved only in the degree
range just specified. It does not assert that
the four variables have an independent joint
law in all degrees.

\subsection{Second and fourth moments}
In the cube class, \eqref{eq:conditional} has
coefficient vector $(-4,-2,-2,-2)$.
The CM contribution has variance $16\cdot2=32$,
and the other three have variance four each.
Thus the variance is 44. In the noncube class,
the coefficient vector is $(2,-2,-2,1)$:
its variance is $4\cdot2+4+4+1=17$.
The first and third moments vanish by
\eqref{eq:functional}.

Only the single-variable fourth powers and
products of two squares survive in the
fourth moment. They give
\begin{align*}
 4224&=4^4\cdot6+3\cdot2^4\cdot2
       +6\{3(32\cdot4)+3(4\cdot4)\},\\
 738&=2^4\cdot6+(2^4+2^4+1)\cdot2
       +6\{8\cdot9+(4\cdot4+4\cdot1+4\cdot1)\}.
\end{align*}
Finally
\[
 \frac{44+2\cdot17}{3}=26,\qquad
 \frac{4224+2\cdot738}{3}=1900.
\]
The uniform bounds following \eqref{eq:conditional}
justify replacing $D_p$ by $Y_p$.
This completes the proof of
Theorem~\ref{thm:moments}.
For the unmultiplied fluctuation
$(A_p-p/27)/\sqrt p$, the second and fourth
moments are $26/729$ and $1900/531441$.

\subsection{The covariance of the original CM pair}
The same lemma applied to \eqref{eq:gamma} gives
\[
 \lim\E\left[
 \begin{pmatrix}v_0^2&v_0v_1\\v_0v_1&v_1^2\end{pmatrix}
 \,\middle|\,3\text{ a cube}\right]
 =\begin{pmatrix}2&2\\2&2\end{pmatrix},
\]
and in the noncube class the matrix is
$\left(\begin{smallmatrix}2&-1\\-1&2\end{smallmatrix}\right)$.
The unconditional covariance is zero,
although the exact twist relations remain.
The arithmetic contrast concerns the cubic
class of \emph{three}. Replacing this stratum by the cube/noncube
class of \emph{two} would change the prime-averaging question.

\section{Exact certificates and reproducibility}\label{sec:certificates}
\subsection{Point orders determine the trace}
For a nonsingular short Weierstrass curve over
$\Fp$, a trace certificate contains a point $P$,
its claimed order $r$ with complete prime
factorization, and a claimed group order $N$.
The verifier checks
\[
 [r]P=O,\qquad [r/q]P\ne O
 \quad\text{for each prime }q\mid r,
\]
and
\begin{equation}\label{eq:hasse-certificate}
 r>2\lfloor2\sqrt p\rfloor,\qquad
 |N-p-1|\le\lfloor2\sqrt p\rfloor,\qquad r\mid N.
\end{equation}
The scalar checks and the complete prime
factorization prove the exact point order.
Lagrange then makes $r$ a divisor of the
true group order. Hasse puts that order
in an integer interval of width
$2\lfloor2\sqrt p\rfloor<r$. There can be
only one multiple of $r$ in it, namely $N$.
Thus $p+1-N$ is the certified trace.

The generator uses baby-step/giant-step in the
Hasse interval and searches at most 256
candidate abscissas. Failure to find a
sufficiently large point order is a refusal,
not a presumed trace. The portable program
uses exact trial primality and a declared
modulus range $p<10^{10}$. This numerical
domain restricts the software, not the
mathematical identities or prime limits.
The verifier checks the point orders without
repeating their search.

\subsection{Three illustrative primes}
The archived exact counts and traces include:
\begin{center}
\begin{tabular}{r rrr r}
\toprule
$p$&$t_1$&$t_2$&$t_3$&$A_p$\\
\midrule
$1\,000\,033$&119&227&$-1069$&$36\,876$\\
$4\,294\,967\,197$&$-34978$&$-101713$&$-19969$&$159\,072\,704$\\
$4\,294\,967\,311$&$-52264$&$-46204$&$-7891$&$159\,072\,584$\\
\bottomrule
\end{tabular}
\end{center}
Their CM triples are, in the same order,
$(-670,1967,-1297)$,
$(43250,85529,-128779)$, and
$(124244,-25963,-98281)$.
All three are in the noncube class of three.
Their exact corrected fluctuation numerators
are $-4355$, $-4145$, and $-7517$.

The first count is also checked against a
direct independent walk through its
$1\,000\,032$ character values. The two
billion-scale examples are obtained from
certificates; their full character words
are not constructed. These examples check
identities and are not a statistical
confirmation sample for the moment limits.

\subsection{What the finite programs establish}
The elliptic verifier checks the formal
V\'elu map, its discriminant, the quotient
change of variables, and small-prime sums.
It rechecks 2625 large point-order
certificates for the fixed-support
applications, as well as the nine
non-CM certificates for the three
quadruplet examples.

The moment verifier checks 207 admissible
primes below 3000, including both classes
of three. It performs 1242 point counts by
two methods, 621 CM--Jacobi comparisons,
621 Fermat-fibre comparisons, and 207
direct four-window checks. It verifies
nine additional CM point-order
certificates at the three illustrative
primes. Arithmetic qualification is
checked at the small primes specified
above. Two separate integer computations,
by multinomials and Laurent constant
terms, reproduce the four moments.
All 70 joint monomials through degree
four and all 315 index cases of
Lemma~\ref{lem:averages} are enumerated.
Altered certificates and requests outside
the stated degree range are rejected.

The 315 cases verify the completeness of
the finite routing table. They do not
prove automorphy or a prime number
theorem. The output moment constants
are exact evaluations of the functional
justified in the proof, not empirical
moments fitted to the 207 primes.

From the repository root, the exact
six-curve oracle can be invoked as
\begin{quote}\ttfamily\small
python3 -B scripts/cubic\_quadruplets.py 4294967197
\end{quote}
It checks all six required curve identities
and certificates and evaluates
\eqref{eq:exact}, without using a fast
Jacobi evaluator. A missing stored
certificate set is reported as such.
The explicit option \texttt{--generate}
requests a bounded new six-curve
certificate search and may refuse.
The two finite certificate replay commands are
\begin{quote}\ttfamily\small
python3 -B reproducibility/cubic\_quadruplets/elliptic/src/verify.py\\
python3 -B reproducibility/cubic\_quadruplets/moments/src/verify.py
\end{quote}
They use standard Python only. These commands check the supplied finite
identities and certificates without a network connection or GPU; they
do not establish the external analytic theorems by computation.

\section{Consequences and scope}
\subsection{A non-Gaussian moment signature}
The ratio of the conditional variances
is $44/17$. Globally the standardized
fourth-moment ratio is
\[
 \frac{1900}{26^2}=\frac{475}{169}.
\]
A centred Gaussian of variance 26 has
fourth moment 2028, not 1900.
Uniform boundedness makes moment
convergence necessary under weak
convergence, so $D_p$ and $Y_p$ cannot
converge in law to that Gaussian.
This conclusion does not establish
existence of a different complete
limiting distribution. Compactness
only assures subsequential limits,
all sharing the four displayed moments.

\subsection{A candidate distribution, not a theorem}
Full joint equidistribution, if proved,
would suggest the mixture
\begin{align*}
 &\frac13\,\mathcal L(8\cos\Theta+2T_1+2T_2+2T_3)\\
 &\qquad+\frac23\,\mathcal L(4\cos\Theta+2T_1+2T_2+T_3),
\end{align*}
where $\Theta$ is uniform and the $T_i$
have semicircle density
$(2\pi)^{-1}\sqrt{4-t^2}$ on $[-2,2]$,
independently in this candidate model.
Symmetry makes the signs immaterial.
Theorem~\ref{thm:moments} verifies its
moments through degree four by a
separate analytic argument. It does
not prove this distributional limit.
For example, the degree-six product
$P_2(x_1)P_2(x_2)P_2(x_3)$ is outside
the finite routing table established
here. Results about one or two curves
are not silently extended to three.

\subsection{Normalized centring and critical maxima}
The bounded correction $b_p$ vanishes
after division by $\sqrt p$.
Vanishing averages of normalized
traces do not imply
\[
 \E(A_p-p/27)=-\E b_p/27+o(1).
\]
That would require substantially
stronger cancellation of the raw
trace sums. No constant-order
raw-centering formula is asserted.

Likewise, a four-window count of
order $p$, with fluctuations of order
$\sqrt p$, does not determine the
rare occurrence count at length
$k\sim\log_3p$ or its probability
of being zero. The fixed-support
application in the appendix retains
exact signs, but makes no uniform
claim about a family whose number
and parameters grow with $k$.
The global maximum problem remains
separate from the theorems here.

\subsection{Manuscript status}
The analytic inputs are the cited literature theorems. Their required
cuspidality, base-change and twist-exclusion conditions are verified for
the fixed representations used here. Finite programs check the arithmetic
identities and the routing of the low-degree calculations, rather than
the external analytic theorems themselves. Specialist review is a
validation task, not an additional mathematical hypothesis or a GRH
assumption.

Artificial-intelligence tools assisted with proof development, computation,
literature retrieval, and English writing and revision. Brice Pouly is the
human author and assumes responsibility for the proof, programs, data,
references and conclusions. Author-controlled checks do not constitute
independent external replication or peer review. The accompanying article
archive contains its numerical certificates and bounded reproduction
commands.

\appendix
\section{The signed support-four sector of a longer start count}
For $k+1<p$ define the predecessor-indexed starts
\[
 J_{p,z}^{(k)}=
 \ind_{\chi(z)\ne\chi(z+1)=\cdots=\chi(z+k)},
 \quad 1\le z\le p-k-1,
 \qquad Z_{p,k}=\sum_zJ_{p,z}^{(k)}.
\]
The neutral prefix has no predecessor in this interval.
At each admissible $z$, subtracting the $(k+1)$-colour
equality indicator from the last-$k$ equality indicator gives
\begin{equation}\label{eq:long-expansion}
 J_{p,z}^{(k)}=3^{-k}
 \sum_{\substack{e\in\{0,1,2\}^{k+1}\\\sum e_j=0\ (3)}}
 (3\ind_{e_0=0}-1)
 \prod_{j:e_j\ne0}\chi(z+j)^{e_j}.
\end{equation}
For a support $S$ of cardinality four, the six
active exponent patterns have two ones and two twos.
They form three conjugate pairs, with common coefficient
$c(S)=-1$ if $0\in S$, and $c(S)=2$ otherwise.
Apply Corollary~\ref{thm:fourpositions} to the three
unoriented partitions of each support.
For integer indices compute $A,\rho$ rationally first
and then reduce modulo $p$.
Grouping $A$ and $-A$ is legitimate because their traces
are equal, as proved above.

Set
\[
 C_k=3\left(2\binom{k}{4}-\binom{k}{3}\right),\qquad
 \mathcal E^{[4]}_{p,k}=
 \sum_{S,\{a,b\}\mid\{c,d\}}c(S)\beta_p(\rho)t_A(p).
\]
The excluded residues are $z=-r$ for $0\le r\le k$.
The active roots already contribute zero; the others
produce the finite correction
\[
 H^{[4]}_{p,k}=
 \sum_{S,\{a,b\}\mid\{c,d\}}c(S)
 \sum_{r\in\{0,\ldots,k\}\setminus S}
 \beta_p\left(\frac{(a-r)(b-r)}{(c-r)(d-r)}\right).
\]
Both the phase weights and this correction use
only small character values of differences.
The exact support-four contribution to $Z_{p,k}$ is
\begin{equation}\label{eq:sector}
 \mathcal L^{[4]}_{p,k}=
 -\frac{2C_k+\mathcal E^{[4]}_{p,k}+H^{[4]}_{p,k}}{3^k}.
\end{equation}
This is a signed identity in the actual count,
not an absolute-value upper bound or a substitute
for the rest of the expansion.

At $k=21$ there are 43890 monomials, 21945
conjugate pairs, and 1169 rational absolute
parameters $|A|$. Zero aggregated weights
can be removed before trace evaluation.
The certified applications are
\begin{center}
\begin{tabular}{r r r r}
\toprule
$p$&$k$&$\mathcal E^{[4]}_{p,k}$&$\mathcal L^{[4]}_{p,k}$\\
\midrule
$1\,000\,033$&13&438273&$-148373/531441$\\
$4\,294\,967\,197$&21&$-108879048$&$36273779/3486784401$\\
$4\,294\,967\,311$&21&130803933&$-43623131/3486784401$\\
\bottomrule
\end{tabular}
\end{center}
The corresponding $H$ values are $-18$, $-6129$,
and 1620. This sector is expressed in units
of the full count $Z_{p,k}$.
The reduction does not estimate all higher
supports collectively, and the four prime-averaged
moments for the six fixed curves provide no
uniform moment theorem for these growing
parameter lists.

\begin{thebibliography}{99}
\bibitem{lehmer}
D. H. Lehmer and E. Lehmer, On runs of residues,
\emph{Proc. Amer. Math. Soc.} 13 (1962), 102--106.
\href{https://doi.org/10.1090/S0002-9939-1962-0138582-6}{doi:10.1090/S0002-9939-1962-0138582-6}.
\bibitem{lehmermachine}
D. H. Lehmer, E. Lehmer, W. H. Mills and J. L. Selfridge,
Machine proof of a theorem on cubic residues,
\emph{Math. Comp.} 16 (1962), 407--415.
\bibitem{graham}
R. L. Graham, On quadruples of consecutive $k$th power residues,
\emph{Proc. Amer. Math. Soc.} 15 (1964), 196--197.
\bibitem{greene}
J. Greene, Hypergeometric functions over finite fields,
\emph{Trans. Amer. Math. Soc.} 301 (1987), 77--101.
\href{https://doi.org/10.1090/S0002-9947-1987-0879564-8}{doi:10.1090/S0002-9947-1987-0879564-8}.
\bibitem{lennon}
C. Lennon, Trace formulas for Hecke operators, Gaussian
hypergeometric functions, and the modularity of a threefold,
\emph{J. Number Theory} 131 (2011), 2320--2351.
\href{https://doi.org/10.1016/j.jnt.2011.05.005}{doi:10.1016/j.jnt.2011.05.005}.
Theorem 1.1 is cited from arXiv:1003.1157v2 (27 December 2010),
\emph{A Trace Formula for Certain Hecke Operators and Gaussian
Hypergeometric Functions}.
\url{https://arxiv.org/abs/1003.1157v2}.
\bibitem{freemansatoh}
D. M. Freeman and T. Satoh, Constructing pairing-friendly
hyperelliptic curves using Weil restriction,
\emph{J. Number Theory} 131 (2011), 959--983.
\href{https://doi.org/10.1016/j.jnt.2010.06.003}{doi:10.1016/j.jnt.2010.06.003}.
\bibitem{guillevicvergnaud}
A. Guillevic and D. Vergnaud, Genus 2 hyperelliptic curve families
with explicit Jacobian order evaluation and pairing-friendly
constructions, \emph{Pairing-Based Cryptography--Pairing 2012},
LNCS 7708 (2013), 234--253.
\href{https://doi.org/10.1007/978-3-642-36334-4_16}{doi:10.1007/978-3-642-36334-4\_16}.
\bibitem{fitesutherland}
F. Fit\'e and A. V. Sutherland, Sato--Tate distributions of twists
of $y^2=x^5-x$ and $y^2=x^6+1$,
\emph{Algebra Number Theory} 8 (2014), 543--585.
\href{https://doi.org/10.2140/ant.2014.8.543}{doi:10.2140/ant.2014.8.543}.
\bibitem{bounds}
B. Pouly, \emph{Succinct Rabung certificates and recurrence-closed lower
bounds for van der Waerden numbers}, published edition v2.0, 2026.
\href{https://doi.org/10.5281/zenodo.23337824}{doi:10.5281/zenodo.23337824}.
\bibitem{atlas}
B. Pouly, \emph{Arithmetic structure in a complete atlas of power-residue
runs below $2^{32}$: Rabung profiles, Kummer densities, norm dependencies,
and local codes}, published edition v1.0, 2026.
\href{https://doi.org/10.5281/zenodo.23339668}{doi:10.5281/zenodo.23339668}.
\bibitem{velu}
J. V\'elu, Isog\'enies entre courbes elliptiques,
\emph{C. R. Acad. Sci. Paris S\'er. A--B} 273 (1971), A238--A241.
\bibitem{newtonthorne}
J. Newton and J. A. Thorne, Symmetric power functoriality
for holomorphic modular forms, II,
\emph{Publ. Math. IH\'ES} 134 (2021), 117--152.
\url{https://arxiv.org/abs/2009.07180}.
\bibitem{ramakrishnan}
D. Ramakrishnan, Modularity of the Rankin--Selberg
$L$-series, and multiplicity one for $\mathrm{SL}(2)$,
\emph{Ann. of Math.} 152 (2000), 45--111.
\url{https://arxiv.org/abs/math/0007203}.
\bibitem{gillespieji}
T. Gillespie and G. Ji, Prime number theorems for
Rankin--Selberg $L$-functions over number fields,
\emph{Sci. China Math.} 54 (2011), 35--46.
\href{https://doi.org/10.1007/s11425-010-4137-x}{doi:10.1007/s11425-010-4137-x}.
Theorem numbering follows arXiv:0910.3660v1 (19 October 2009),
\emph{A Prime number theorem for Rankin--Selberg $L$-functions
over number fields}. \url{https://arxiv.org/abs/0910.3660v1}.
\bibitem{chenshen}
Q. Chen and E. Shen, Effective Sato--Tate distributions
for surfaces arising from products of elliptic curves,
\emph{Res. Math. Sci.} 11 (2024), article 69.
\href{https://doi.org/10.1007/s40687-024-00470-x}{doi:10.1007/s40687-024-00470-x}.
Section 3.3 and Theorem 3.6 refer to arXiv:2309.08848v1
(16 September 2023), whose title starts \emph{On Effective
Sato--Tate distributions}. \url{https://arxiv.org/abs/2309.08848v1}.
\bibitem{sutherland13}
A. V. Sutherland, \emph{18.783 Elliptic Curves},
MIT lecture notes 13 (2023).
\url{https://math.mit.edu/classes/18.783/2023/LectureNotes13.pdf}.
\bibitem{sutherland21}
A. V. Sutherland, \emph{18.783 Elliptic Curves},
MIT lecture notes 21 (2023), Section 21.5.
\url{https://math.mit.edu/classes/18.783/2023/LectureNotes21.pdf}.
\end{thebibliography}
\end{document}
